\chapter*{Thông tin đề tài và nhóm thực hiện}
\addcontentsline{toc}{chapter}{Thông tin đề tài và nhóm thực hiện}
\begin{center}
{\Large\bfseries BÁO CÁO TỔNG KẾT BTL03}\\[0.5em]
{\large Chuyển đổi mã nguồn Deep Learning từ Keras sang PyTorch}\\[1em]
Snapshot: T0212 / commit \texttt{33dcd74599e86404bbc8128c8057afba82810093}\\
Ngày chốt evidence: 2026-11-23\\[0.5em]
Alias thành viên có bằng chứng: Duy-DS; nguyenthanhdat180905-cloud; nguyen-cong-tri; TonOfWeight072501; CHIVY2005\\
Họ tên, MSSV, giảng viên và tỷ lệ đóng góp: [CẦN NHÓM XÁC NHẬN; tổng tỷ lệ phải bằng 100\%]
\end{center}

\chapter{Tổng quan dự án}
\section{Bối cảnh và mục tiêu}
BTL03 chuyển đổi các mô hình tham chiếu Keras sang PyTorch theo cấu trúc hướng đối tượng, đồng thời kiểm tra layout, loss, vòng lặp huấn luyện và quản lý thực nghiệm. Báo cáo này là snapshot tiến độ, không phải bản nghiệm thu cuối.

\section{Phạm vi}
Mười cấu hình thuộc năm nhóm: MLP; CNN+MLP; CNN+LSTM; utilities/callbacks; SegNet và FCN score-map. Hai cấu hình utilities dùng lại Gender/Emotion, do đó mười cấu hình không phải mười kiến trúc độc lập.

\section{Trạng thái tại T0212}
Report T0212 ghi 10/10 synthetic one-epoch cases và parameter-count audit đạt, nhưng logger bị chặn bởi protobuf/TensorBoard; full suite có bốn lỗi blocker. Trạng thái tổng thể là \textbf{BLOCKED}, không đủ cơ sở tuyên bố nghiệm thu.

\section{Tổ chức báo cáo}
Chương 2 trình bày khác biệt framework; chương 3 mô tả kiến trúc; chương 4 kiểm toán/thực nghiệm; chương 5 nêu giới hạn và hướng phát triển.